\chapter{Gram–Schmidt and Orthogonal Projection}
\label{ch:i14-gram-schmidt-and-orthogonal-projection}

Gram–Schmidt converts an arbitrary independent list into an orthonormal list spanning the same nested subspaces. Orthogonal projection then gives canonical decompositions and solves least-squares problems. This chapter connects exact geometry with the numerical algorithms QR factorization and normal equations.

\section*{Learning goals}
After this chapter, the reader should be able to:
\begin{itemize}
\item perform the Gram--Schmidt process on an independent list and verify that successive spans are preserved;
\item compute orthogonal projections onto lines and higher-dimensional subspaces from an orthonormal basis;
\item prove the best-approximation property by decomposing the error into orthogonal components;
\item construct a projection matrix and verify that it is idempotent and self-adjoint;
\item derive a reduced QR factorization from Gram--Schmidt;
\item solve a least-squares problem using projection, normal equations, or QR and compare the numerical consequences;
\end{itemize}
\section*{Conceptual roadmap}
The chapter is organized around a repeated pattern:
\[
\boxed{\text{definition}\;\longrightarrow\;\text{structural theorem}\;\longrightarrow\;\text{coordinate calculation}\;\longrightarrow\;\text{application}.}
\]
Whenever possible, we separate the invariant mathematical object from a chosen coordinate representation.

\section{Orthonormal lists}
A list is orthonormal when its vectors have norm one and are pairwise orthogonal. Coordinates in an orthonormal basis are inner products.

\section{Gram–Schmidt process}
Successively subtract from each vector its projections onto the previously constructed orthonormal vectors, then normalize.


% BEGIN VOL01-EXPANSION I14-example-01
\begin{example}[Gram--Schmidt with a verification]\label{ex:i14-expand-01}
Start with
\[
v_1=(1,1,0),\qquad v_2=(1,0,1).
\]
Set \(e_1=v_1/\sqrt2\). Then
\[
u_2=v_2-\langle v_2,e_1\rangle e_1
=(1,0,1)-\frac12(1,1,0)
=\left(\frac12,-\frac12,1\right).
\]
Since \(\|u_2\|=\sqrt{3/2}\), set \(e_2=u_2/\sqrt{3/2}\).
A direct dot-product check gives \(\langle e_1,e_2\rangle=0\), and both have
norm \(1\). The construction also preserves the span because
\(u_2=v_2-\frac12v_1\).
\end{example}
% END VOL01-EXPANSION I14-example-01
\section{Projection onto a subspace}
For an orthonormal basis $e_1,\ldots,e_r$ of $U$, $P_Uv=\sum_i\langle v,e_i\rangle e_i$.


% BEGIN VOL01-EXPANSION I14-example-02
\begin{example}[Projection matrix onto a plane]\label{ex:i14-expand-02}
Let \(U\subset\mathbb R^3\) have orthonormal basis
\[
q_1=(1,0,0),\qquad q_2=(0,1/\sqrt2,1/\sqrt2).
\]
With \(Q=[q_1\ q_2]\), the orthogonal projector is
\[
P=QQ^T.
\]
For \(v=(2,3,1)\),
\[
P v=2q_1+\frac{4}{\sqrt2}q_2=(2,2,2).
\]
The residual \((0,1,-1)\) has zero inner product with both \(q_1,q_2\),
checking that it lies in \(U^\perp\). Also \(P^2=P\), as every orthogonal
projector must satisfy.
\end{example}
% END VOL01-EXPANSION I14-example-02
\section{Best approximation}
$P_Uv$ is the unique vector in $U$ minimizing $\|v-u\|$.

\section{QR factorization}
Applying Gram–Schmidt to independent columns of $A$ produces $A=QR$, with $Q$ having orthonormal columns and $R$ upper triangular.

\section{Least squares}
For full-column-rank $A$, the least-squares solution satisfies $A^*Ax=A^*b$ and can be computed more stably using QR.


% BEGIN VOL01-EXPANSION I14-example-03
\begin{example}[Least squares from QR and residual geometry]\label{ex:i14-expand-03}
Fit \(ax\approx b\) with
\[
A=\begin{pmatrix}1\\1\\1\end{pmatrix},\qquad
b=\begin{pmatrix}1\\2\\4\end{pmatrix}.
\]
The unit column is \(q=A/\sqrt3\), so the least-squares fitted vector is
\[
qq^Tb=\frac{7}{3}\begin{pmatrix}1\\1\\1\end{pmatrix}.
\]
Thus \(x=7/3\). The residual is
\[
r=b-Ax=\left(-\frac43,-\frac13,\frac53\right)^T,
\]
and \(A^Tr=0\), which verifies the normal-equation condition. The method
therefore agrees with the geometric statement that the residual is orthogonal
to the column space.
\end{example}
% END VOL01-EXPANSION I14-example-03
\section{Core structural results}

\begin{theorem}[Gram–Schmidt theorem]\label{thm:i14-01}
Every finite independent list in an inner-product space can be replaced by an orthonormal list with the same successive spans.
\begin{proof}
At step $k$, subtract the projection of $v_k$ onto the span already constructed. Independence guarantees the remainder is nonzero; normalization finishes the step.
\end{proof}
\end{theorem}

\begin{theorem}[Projection theorem]\label{thm:i14-02}
For finite-dimensional $U$, every $v$ decomposes uniquely as $v=P_Uv+(v-P_Uv)$ with the first term in $U$ and the second in $U^\perp$.
\begin{proof}
Use an orthonormal basis of $U$ to define $P_Uv$. Direct calculation shows the residual is orthogonal to every basis vector and hence to $U$.
\end{proof}
\end{theorem}

\section{Worked examples}

\begin{example}[Two planar vectors]\label{ex:i14-worked-01}
Starting from $(1,0)$ and $(1,1)$, Gram–Schmidt returns $(1,0)$ and $(0,1)$.
\end{example}

\begin{example}[Projection onto a line]\label{ex:i14-worked-02}
Projection onto the line spanned by unit $u$ is $P(v)=\langle v,u\rangle u$.
\end{example}

\begin{example}[Least squares line fitting]\label{ex:i14-worked-03}
The columns $1$ and $t$ form a design matrix; projection of data onto their column space gives the best affine fit.
\end{example}

\section{Diagnostic checklist}
Before moving on, it is useful to distinguish four levels of understanding:
\begin{enumerate}
\item recognize the definitions in unfamiliar examples;
\item prove the core structural statements without coordinates when possible;
\item translate the same statement into a matrix or coordinate computation;
\item know which hypotheses are essential and produce a counterexample when one is removed.
\end{enumerate}

% BEGIN VOLUME01 SOLVED DOSSIERS
\section{Solved dossiers}
These dossiers are longer worked problems. They complement the shorter exercise--hint--solution triads by combining several ideas in one argument.

\begin{problem}[Gram–Schmidt by hand]\label{prob:i14-dossier-01}
Apply Gram–Schmidt to $v_1=(1,1,0)$ and $v_2=(1,0,1)$.
\end{problem}
\begin{solution}
Set $e_1=v_1/\sqrt2$. The projection coefficient is $\langle v_2,e_1\rangle=1/\sqrt2$, so $u_2=v_2-\frac12(1,1,0)=(1/2,-1/2,1)$. Its norm is $\sqrt{3/2}$, hence $e_2=u_2/\sqrt{3/2}$.
\end{solution}

\begin{problem}[Projection onto a line]\label{prob:i14-dossier-02}
Project $v=(2,3)$ onto the line spanned by $u=(1,1)$.
\end{problem}
\begin{solution}
Using the unit vector $e=u/\sqrt2$, $P(v)=\langle v,e\rangle e=(5/\sqrt2)(1/\sqrt2,1/\sqrt2)=(5/2,5/2)$.
\end{solution}

\begin{problem}[Projection matrix]\label{prob:i14-dossier-03}
Find the projection matrix onto the line spanned by a unit vector $u$.
\end{problem}
\begin{solution}
$P=uu^T$ in the real case. Indeed $Pv=u(u^Tv)$ is exactly the scalar projection coefficient times $u$.
\end{solution}

\begin{problem}[Best approximation]\label{prob:i14-dossier-04}
Why is $P_Uv$ the unique closest vector in $U$ to $v$?
\end{problem}
\begin{solution}
For $u\in U$, write $v-u=(v-P_Uv)+(P_Uv-u)$ with orthogonal summands. Pythagoras gives $\|v-u\|^2=\|v-P_Uv\|^2+\|P_Uv-u\|^2$, minimized uniquely at $u=P_Uv$.
\end{solution}

\begin{problem}[QR factorization]\label{prob:i14-dossier-05}
What does Gram–Schmidt do to the columns of a full-rank matrix $A$?
\end{problem}
\begin{solution}
It produces orthonormal columns $Q$ and expresses each original column as a combination of the current and previous orthonormal columns. The coefficient matrix is upper triangular $R$, giving $A=QR$.
\end{solution}

\begin{problem}[Normal equations]\label{prob:i14-dossier-06}
Derive the least-squares normal equations.
\end{problem}
\begin{solution}
At the minimizer, the residual $b-Ax$ is orthogonal to the column space of $A$. Therefore $A^*(b-Ax)=0$, which rearranges to $A^*Ax=A^*b$.
\end{solution}

\begin{problem}[Why QR is stabler]\label{prob:i14-dossier-07}
Why is solving least squares through QR often preferable to normal equations?
\end{problem}
\begin{solution}
Forming $A^*A$ squares the spectral condition number, magnifying sensitivity. Orthogonal transformations used in QR have condition number $1$ and preserve norms, so they avoid that amplification.
\end{solution}

\begin{problem}[Projection idempotence]\label{prob:i14-dossier-08}
Prove an orthogonal projection satisfies $P^2=P$.
\end{problem}
\begin{solution}
After one projection, $Pv$ lies in the target subspace. Projecting a vector already in that subspace leaves it fixed, hence $P(Pv)=Pv$.
\end{solution}

\begin{problem}[Complementary projection]\label{prob:i14-dossier-09}
Show $I-P_U$ is the orthogonal projection onto $U^\perp$.
\end{problem}
\begin{solution}
The orthogonal decomposition is $v=P_Uv+(v-P_Uv)$. The second term lies in $U^\perp$ and equals $(I-P_U)v$. Vectors already in $U^\perp$ are fixed.
\end{solution}

\begin{problem}[Orthonormal coordinates]\label{prob:i14-dossier-10}
If $e_1,\dots,e_n$ is orthonormal, why is the coefficient of $e_j$ in $v$ equal to $\langle v,e_j\rangle$?
\end{problem}
\begin{solution}
Write $v=\sum a_ie_i$ and inner-product with $e_j$. Orthonormality kills all terms except $a_j\langle e_j,e_j\rangle=a_j$.
\end{solution}

\begin{problem}[Rank-deficient least squares]\label{prob:i14-dossier-11}
What changes if the columns of $A$ are dependent?
\end{problem}
\begin{solution}
The least-squares projection onto the column space still exists uniquely, but the coefficient vector $x$ need not be unique because nonzero kernel vectors can be added without changing $Ax$. The pseudoinverse selects the minimum-norm solution.
\end{solution}

\begin{problem}[Projection synthesis]\label{prob:i14-dossier-12}
Connect Gram–Schmidt, QR, and least squares in one sentence.
\end{problem}
\begin{solution}
Gram–Schmidt constructs an orthonormal basis of the column space, QR records that construction algebraically, and least squares projects data onto that same column space using the orthonormal coordinates.
\end{solution}

% END VOLUME01 SOLVED DOSSIERS

\section{Exercises with complete solutions}

\begin{exercise}\label{exr:i14-01}
Apply Gram–Schmidt to $(1,1)$ and $(1,-1)$.
\end{exercise}
\begin{hint}
At the next Gram--Schmidt step, subtract projections onto the orthonormal vectors already obtained before normalizing the residual. At each Gram--Schmidt step, subtract only the components along the orthonormal vectors already constructed.
\end{hint}
\begin{solution}
They are already orthogonal; normalize each by $\sqrt2$.
\end{solution}

\begin{exercise}\label{exr:i14-02}
Project $(2,3)$ onto the $x$-axis.
\end{exercise}
\begin{hint}
For projection onto a line spanned by unit \(u\), start with \(P(v)=\langle v,u\rangle u\).
\end{hint}
\begin{solution}
The projection is $(2,0)$.
\end{solution}

\begin{exercise}\label{exr:i14-03}
Why is the projection residual orthogonal to the subspace?
\end{exercise}
\begin{hint}
Write the relevant inner products explicitly and use bilinearity or sesquilinearity before simplifying norms or angles. Compute the residual \(v-P_Uv\) and test its inner product with every basis vector of \(U\).
\end{hint}
\begin{solution}
Subtracting all orthonormal-basis components makes its inner product with each basis vector zero.
\end{solution}

\begin{exercise}\label{exr:i14-04}
Prove the best-approximation property.
\end{exercise}
\begin{hint}
Decompose the error into the projection residual and a vector inside the subspace; orthogonality turns the norm square into a Pythagorean sum. Decompose \(v-u=(v-P_Uv)+(P_Uv-u)\); then use orthogonality and Pythagoras.
\end{hint}
\begin{solution}
For $u\in U$, write $v-u=(v-P_Uv)+(P_Uv-u)$; the two terms are orthogonal, so Pythagoras is minimized exactly when $u=P_Uv$.
\end{solution}

\begin{exercise}\label{exr:i14-05}
What are the dimensions of $Q$ and $R$ for an $m\times n$ matrix of full column rank?
\end{exercise}
\begin{hint}
Choose a basis adapted to the invariant flag; the vanishing entries below the diagonal encode preservation of the successive subspaces. In reduced QR for an \(m\times n\) full-column-rank matrix, count the orthonormal columns of \(Q\) before deciding dimensions.
\end{hint}
\begin{solution}
In reduced QR, $Q$ is $m\times n$ with orthonormal columns and $R$ is invertible $n\times n$ upper triangular.
\end{solution}

\begin{exercise}\label{exr:i14-06}
Why can normal equations be less stable than QR?
\end{exercise}
\begin{hint}
Use that the least-squares residual must be orthogonal to the column space, so apply \(A^*\) to \(b-Ax\). Compare \(\kappa(A^*A)\) with \(\kappa(A)\); forming the normal equations squares the condition number.
\end{hint}
\begin{solution}
Forming $A^*A$ squares the condition number, amplifying numerical sensitivity.
\end{solution}

\begin{exercise}\label{exr:i14-07}
Show an orthogonal projection satisfies $P^2=P$.
\end{exercise}
\begin{hint}
After the first projection, the vector already lies in the target subspace; project that vector one more time. After one projection the vector is already in \(U\); apply \(P\) once more.
\end{hint}
\begin{solution}
Once a vector lies in $U$, projecting again leaves it unchanged; hence $P(Pv)=Pv$.
\end{solution}

\begin{exercise}\label{exr:i14-08}
Show $I-P_U$ projects onto $U^\perp$.
\end{exercise}
\begin{hint}
Use the orthogonal decomposition \(v=P_Uv+(v-P_Uv)\) and identify where the residual lies. Use \(v=P_Uv+(v-P_Uv)\) and identify the second term as the \(U^\perp\) component.
\end{hint}
\begin{solution}
The decomposition theorem gives $v-P_Uv\in U^\perp$, and vectors already in $U^\perp$ are fixed by $I-P_U$.
\end{solution}


% BEGIN VOL01-EXPANSION I14-exercises-01
\section{Graded supplementary exercises}
These exercises extend the preserved foundational set and are graded by mathematical role.

\subsection*{Standard computations}

\begin{exercise}[Gram--Schmidt pair]\label{exr:i14-expand-01}
Apply Gram--Schmidt to \((1,1,0)\) and \((1,0,1)\).
\end{exercise}
\begin{hint}
Normalize the first vector, subtract its projection from the second, then normalize.
\end{hint}
\begin{solution}
One obtains \(e_1=(1,1,0)/\sqrt2\) and \(e_2=(1,-1,2)/\sqrt6\).
\end{solution}

\begin{exercise}[Projection onto a line]\label{exr:i14-expand-02}
Project \(v=(3,1)\) onto \(\operatorname{span}(1,2)\).
\end{exercise}
\begin{hint}
Use \(P_Uv=\langle v,u\rangle u/\langle u,u\rangle\).
\end{hint}
\begin{solution}
The coefficient is \(5/5=1\), so the projection is \((1,2)\).
\end{solution}

\begin{exercise}[Projection matrix]\label{exr:i14-expand-03}
Find the matrix projecting \(\mathbb R^2\) onto \(\operatorname{span}(1,1)\).
\end{exercise}
\begin{hint}
Use the unit vector \(u=(1,1)/\sqrt2\) and \(P=uu^T\).
\end{hint}
\begin{solution}
\(P=\frac12\begin{pmatrix}1&1\\1&1\end{pmatrix}\).
\end{solution}

\begin{exercise}[Reduced QR]\label{exr:i14-expand-04}
Find a reduced QR factorization of \(A=\begin{pmatrix}1&1\\0&1\\0&0\end{pmatrix}\).
\end{exercise}
\begin{hint}
Apply Gram--Schmidt to the two columns.
\end{hint}
\begin{solution}
Here \(q_1=e_1\), \(u_2=(0,1,0)\), \(q_2=e_2\); \(Q=[e_1,e_2]\), \(R=\begin{pmatrix}1&1\\0&1\end{pmatrix}\).
\end{solution}

\begin{exercise}[Normal equation]\label{exr:i14-expand-05}
Solve the least-squares problem for \(A=(1,1,1)^T\), \(b=(1,2,4)^T\).
\end{exercise}
\begin{hint}
Solve \(A^TAx=A^Tb\).
\end{hint}
\begin{solution}
\(3x=7\), so \(x=7/3\).
\end{solution}

\subsection*{Proofs}

\begin{exercise}[Gram--Schmidt span preservation]\label{exr:i14-expand-06}
Prove each Gram--Schmidt step preserves the span of the vectors processed so far.
\end{exercise}
\begin{hint}
The new residual differs from the new original vector by a combination of earlier vectors.
\end{hint}
\begin{solution}
Each residual lies in the old span, while the original vector is recovered from the residual plus earlier projection terms; hence the successive spans agree.
\end{solution}

\begin{exercise}[Projection idempotence]\label{exr:i14-expand-07}
Prove an orthogonal projection satisfies \(P^2=P\).
\end{exercise}
\begin{hint}
After one projection the vector already lies in the target subspace.
\end{hint}
\begin{solution}
For every \(v\), \(Pv\in U\), and projection fixes every vector of \(U\); hence \(P(Pv)=Pv\).
\end{solution}

\begin{exercise}[Projection self-adjointness]\label{exr:i14-expand-08}
Prove \(P_U^*=P_U\).
\end{exercise}
\begin{hint}
Use an orthonormal basis \(q_i\) and the formula \(P_Uv=\sum\langle v,q_i\rangle q_i\).
\end{hint}
\begin{solution}
Expanding both \(\langle P_Uv,w\rangle\) and \(\langle v,P_Uw\rangle\) gives the same sum.
\end{solution}

\begin{exercise}[Best approximation]\label{exr:i14-expand-09}
Prove \(P_Uv\) uniquely minimizes \(\|v-u\|\) over \(u\in U\).
\end{exercise}
\begin{hint}
Decompose \(v-u=(v-P_Uv)+(P_Uv-u)\).
\end{hint}
\begin{solution}
The two terms are orthogonal, so Pythagoras gives a fixed residual term plus \(\|P_Uv-u\|^2\), minimized only at \(u=P_Uv\).
\end{solution}

\subsection*{Counterexamples and hypothesis testing}

\begin{exercise}[Dependent Gram--Schmidt input]\label{exr:i14-expand-10}
What fails if Gram--Schmidt is applied to a dependent list?
\end{exercise}
\begin{hint}
At some step the residual can become zero.
\end{hint}
\begin{solution}
For \(v_2=2v_1\), subtracting the projection of \(v_2\) onto \(v_1\) leaves zero, which cannot be normalized.
\end{solution}

\begin{exercise}[Normal equations conditioning]\label{exr:i14-expand-11}
Why can forming \(A^TA\) be a poor numerical choice?
\end{exercise}
\begin{hint}
Compare singular values of \(A\) and \(A^TA\).
\end{hint}
\begin{solution}
The condition number is squared: \(\kappa_2(A^TA)=\kappa_2(A)^2\) for full column rank.
\end{solution}

\begin{exercise}[Oblique idempotent]\label{exr:i14-expand-12}
Give an idempotent matrix that is not an orthogonal projector.
\end{exercise}
\begin{hint}
Try an upper-triangular rank-one projection that is not symmetric.
\end{hint}
\begin{solution}
\(P=\begin{pmatrix}1&1\\0&0\end{pmatrix}\) satisfies \(P^2=P\) but \(P^T\ne P\), so it is not an orthogonal projection.
\end{solution}

\subsection*{Applications and investigations}

\begin{exercise}[Least-squares line]\label{exr:i14-expand-13}
Explain how fitting \(y=a+bt\) to data becomes a projection problem.
\end{exercise}
\begin{hint}
Use design columns \(1\) and \(t\).
\end{hint}
\begin{solution}
The fitted data vector is the orthogonal projection of the observation vector onto the column space of the design matrix.
\end{solution}

\begin{exercise}[QR stability]\label{exr:i14-expand-14}
Why are orthogonal transformations used in numerical least squares?
\end{exercise}
\begin{hint}
Recall their Euclidean condition number and norm preservation.
\end{hint}
\begin{solution}
Orthogonal transformations have condition number \(1\), so they reorganize the problem without amplifying Euclidean perturbations.
\end{solution}

\subsection*{Challenge problems}

\begin{exercise}[Projection characterization]\label{exr:i14-expand-15}
Prove a real matrix \(P\) is an orthogonal projector iff \(P^2=P\) and \(P^T=P\).
\end{exercise}
\begin{hint}
Use the image/kernel decomposition and show \(\ker P=(\operatorname{im}P)^\perp\).
\end{hint}
\begin{solution}
Idempotence gives \(V=\operatorname{im}P\oplus\ker P\); symmetry gives orthogonality of these spaces, so \(P\) is the orthogonal projection onto its image. The converse is standard.
\end{solution}

\begin{exercise}[Least-squares solution set]\label{exr:i14-expand-16}
If \(A\) is rank deficient and \(x_0\) is one least-squares solution, describe all coefficient vectors giving the same fitted vector.
\end{exercise}
\begin{hint}
Add vectors in \(\ker A\).
\end{hint}
\begin{solution}
All are \(x_0+z\) with \(z\in\ker A\); they produce the same \(Ax\). The pseudoinverse selects the minimum-norm member.
\end{solution}
% END VOL01-EXPANSION I14-exercises-01
\section*{Chapter summary}
Gram–Schmidt converts an arbitrary independent list into an orthonormal list spanning the same nested subspaces. Orthogonal projection then gives canonical decompositions and solves least-squares problems. This chapter connects exact geometry with the numerical algorithms QR factorization and normal equations.
The important habit is to move fluently between invariant statements and concrete coordinates while remembering which conclusions depend on finite dimensionality, the scalar field, or the inner product.
